\guidechapter{From long multiplication to convolution}{A complete small example.}

Write an integer in radix $B$ as $a_0+a_1B+\cdots+a_{q-1}B^{q-1}$.
The digits are listed from least significant to most significant. For the
explanation take $B=10$; the paper uses a power of two so that digit splitting
is a binary operation. Introduce a formal variable $x$ and form
$A(x)=\sum_i a_i x^i$. Then the original integer is exactly $A(B)$.

For two inputs, the coefficients of the polynomial product satisfy
\[
 A(x)B(x)=\sum_k c_k x^k,\qquad c_k=\sum_{i+j=k}a_i b_j.
\]
The second input polynomial is denoted $B(x)$ here; the unparenthesized $B$
still denotes the radix. The sum defining $c_k$ is \emph{convolution}.
It groups digit products that have the same place value. Unlike digits,
the coefficients $c_k$ can exceed the radix. Carry propagation repairs that
after the coefficients have been computed.

\explfigure{01-convolution}{One multiplication, three representations.}{
The coefficient list $[12,11,2]$ belongs to powers $x^0,x^1,x^2$.
The carry is passed from lower to higher positions. The result is $322$.}

For $23\cdot14$, the coefficient of $x$ is $3\cdot1+2\cdot4=11$.
At position zero, $12$ writes digit $2$ and carries $1$. Position one
receives $11+1=12$, again writing $2$ and carrying $1$. Position two
receives $2+1=3$. This is ordinary multiplication reorganized, with no
approximation so far.

\paragraph{Why transforms enter.}
For a suitable transform $F$, convolution can be recovered from products of
corresponding transform entries. Evaluating two polynomials at many roots of
unity, multiplying those values, and transforming back replaces many
cross-products by structured additions and pointwise products. The benefit
depends on computing the transforms and arranging their inputs cheaply.
The paper's transforms are normalized, so known scale factors must be put back;
the slogan ``transform, multiply, inverse transform'' hides those factors.

\paragraph{Why padding is necessary.}
A cyclic convolution treats positions modulo a length $S$. A contribution at
position $S$ then wraps to position zero. That is useful algebraically but
would corrupt this integer product. For degree at most $2q-2$, choosing
$S>2q-2$ and padding unused entries with zero prevents nonzero wraparound.
The original assembly proves its chosen prime-axis product has enough capacity.
It also bounds each coefficient before choosing its storage precision.

\takeaway{The hard task is now to compute all convolution coefficients
cheaply and accurately enough to recover the integers, then carry once.}
\sourceanchor{00-introduction.tex; 08-assembly.tex, eq:capacity}
